% 概率与博弈知识补充；由 mathematics.tex 引入，不单独建立 chapter。
\section{期望选型：计数拆指标，等待拆尾和}\label{knowledge-expectation}\index{indicator variable}\index{tail sum}
\textbf{先确定随机对象。} 均匀随机排列、逐步等概率转移和等概率选择最终方案是三种不同模型。
例如从起点等概率走向 A 或 B，A 只有一个终点，B 有两个等概率终点：
三个终点概率为 $1/2,1/4,1/4$，不是各 $1/3$。不能直接用路径条数代替路径概率。

\textbf{计数型答案：给每个对象一个指示变量。}
若 $X=\sum_{i=1}^m I_i$，其中 $I_i=[\text{事件 }A_i\text{ 发生}]$，则
\[
\mathbb E X=\sum_{i=1}^m\Pr(A_i).
\]
有限和的期望线性性不要求独立；一般实值变量应保证期望存在，避免相减两个无穷大。
随机排列的不动点数为 $X=\sum_i[\pi(i)=i]$，每个位置固定的概率为 $1/n$，所以
$n\ge1$ 时 $\mathbb E X=1$。这比先求全部不动点个数的分布容易。
$n=2$ 时两个指示变量完全相同，恰好说明它们并不独立。

线性性不能推广成 $\mathbb E(XY)=\mathbb E X\mathbb E Y$。
令 $X=Y$ 为一次公平硬币正面的指示变量，则左边为 $1/2$、右边为 $1/4$。
方差应写成
\[
\operatorname{Var}\!\left(\sum_iX_i\right)
=\sum_i\operatorname{Var}(X_i)+2\sum_{i<j}\operatorname{Cov}(X_i,X_j).
\]
两两独立足以消去协方差；仅仅知道各自成功概率不足以计算方差。

\textbf{等待时间或最大值：先问是否超过阈值。}
非负整数随机变量逐点满足 $X=\sum_{k\ge1}[X\ge k]$，因此
\[
\mathbb E X=\sum_{k\ge1}\Pr(X\ge k)
=\sum_{k\ge0}\Pr(X>k).
\]
非负项允许交换求和，结果也可能是 $+\infty$。两种写法的起点相差一。
例如两个独立公平六面骰子的最大值 $M$ 满足
$\Pr(M<k)=((k-1)/6)^2$，故
\[
\mathbb E M=\sum_{k=1}^{6}\left(1-\frac{(k-1)^2}{36}\right)
=6-\frac{55}{36}=\frac{161}{36}.
\]
这里乘出平方用到了两个骰子独立；若两个骰子总是同点，期望就变成 $7/2$。

\textbf{覆盖过程：按完成一件新事的阶段拆分。}
每次独立均匀抽取 $n$ 种券，已有 $i$ 种时，下一张为新种类的概率为 $(n-i)/n$。
含成功那次在内的等待时间期望是 $n/(n-i)$，收齐的期望为
$n\sum_{j=1}^n1/j$。每次抽样方式必须保持不变；非均匀抽券时，“已见种类数”通常不是充分状态，
因为尚缺哪几种会改变转移概率，不能直接套调和数。

\section{条件期望：先定义状态，再解自环}\label{knowledge-probability-dp}\index{conditional expectation}
状态必须保存决定未来分布所需的信息。若到达终点后不再花费，设终点 $t$ 的 $E_t=0$。
一次转移从 $i$ 到 $j$ 的概率为 $p_{ij}$，本次费用为 $c_{ij}$，则第一步分析给出
\[
E_i=\sum_jp_{ij}(c_{ij}+E_j)=b_i+\sum_jp_{ij}E_j,
\qquad b_i=\sum_jp_{ij}c_{ij}.
\]
每步统一花费 1 时才有 $b_i=1$；已经是终点时不能再多加一次。
若目标是命中某个好终点的概率，则好终点边界为 1、坏终点为 0，非终点方程为
$h_i=\sum_jp_{ij}h_j$，没有常数 1。
若过程可能永不终止，命中概率是边界问题的最小非负解，不能把任意代数解当成概率。

\textbf{只有自环时可以局部移项。} 若 $p_{ii}<1$，则
\[
E_i=\frac{b_i+\sum_{j\ne i}p_{ij}E_j}{1-p_{ii}}.
\]
其余边若指向已算出的状态，就按依赖顺序做 DP。不能在存在 $i\to j\to i$ 时只去掉自环后仍当 DAG。
每次独立以概率 $q>0$ 成功，失败回到原状，每次花费 1，有
$E=1+(1-q)E$，故 $E=1/q$。当 $q=0$ 时永不成功，不能继续做除法。
若统计“成功前失败次数”，当前成功不计费用，结果才是 $(1-q)/q$。

\textbf{小例：等到连续两个正面。} 公平硬币独立抛掷，状态 0 表示当前没有末尾正面，
状态 1 表示末尾恰有一个正面，状态 2 表示已完成。第一步方程为
\[
E_0=1+\tfrac12E_0+\tfrac12E_1,\qquad
E_1=1+\tfrac12E_0,\qquad E_2=0.
\]
解得 $E_0=6,E_1=4$。错误地把每两次抛掷作为成功率 $1/4$ 的独立尝试会得到 8，
它只检查不重叠的分组，漏掉跨组的连续正面。反过来把所有相邻两次都当独立试验也不成立。

\textbf{有环时检查吸收，再选择消元。}
只保留起点可达、尚未吸收的状态，记它们之间的转移矩阵为 $Q$。
有限期望向量满足 $(I-Q)E=b$；稠密直接消元的时间为 $O(s^3)$、空间为 $O(s^2)$，
其中 $s$ 是非终点状态数。DAG、链、带状矩阵或可按强连通分量分块的问题应先利用结构，
不必一律构建全局稠密矩阵。概率极小时实数消元可能病态；残差小也不保证答案相对误差小。
实数接口见\hyperref[compact-GaussReal]{GaussReal}（第~\pageref{compact-GaussReal}~页，\ref{compact-GaussReal}~节）；
其秩与一致性按显式阈值判定，不能代替精确可解性或下述吸收条件的证明。

对有限状态、固定转移概率的链，若每个可达非终点都有一条正概率路径通向吸收态，
则存在统一的 $L$ 和 $\varepsilon>0$，使任何非终点在未来 $L$ 步内吸收的概率至少为 $\varepsilon$。
由此 $\Pr(T>kL)\le(1-\varepsilon)^k$，所以几乎必然吸收且 $\mathbb E T\le L/\varepsilon<\infty$。
等价的图判据是：可达区域中没有封闭的非吸收强连通分量。
这里的边只取\emph{正概率}转移，且每步费用有界时累计费用也具有有限绝对期望。

\textbf{两类反例不能混淆。} 以概率 $1/2$ 立即终止、以概率 $1/2$ 进入永久自环，
有正概率永不吸收，停止时间期望为无穷大。
另取正整数随机变量 $T$，令 $\Pr(T=k)=1/(k(k+1))$。
概率和为 1，所以 $T$ 几乎必然有限；但 $\Pr(T\ge k)=1/k$，故 $\mathbb E T=\infty$。
第二例说明无限状态或无限支撑下，“几乎必然终止”不足以推出有限期望，不能丢掉上述有限状态条件。

\section{模概率：可逆的不只是输入分母}\label{knowledge-mod-probability}
竞赛中的“概率对素数 $p$ 取模”不是一个介于 0 与 1 的数，而是有理数的代数表示。
既约分数 $a/b$ 只有在 $p\nmid b$ 时，才能映射为 $a\,b^{-1}\pmod p$。
模值不能比较大小、判定真实概率是否很小，也不能由模值为零断言真实概率为零。
例如概率 $p/(p+1)>0$，其模值为零。

\textbf{直接转模方程的充分条件：}
每个转移概率和费用的分母均可逆，且 $I-Q$ 在模 $p$ 下可逆。
实数下可逆不保证模下可逆；自环移项中的 $1-p_{ii}$ 也必须非零。
例如每次以 $p/(p+1)$ 成功，期望为 $(p+1)/p$；
原始转移概率的分母 $p+1$ 可逆，但期望的分母 $p$ 不可逆。
模方程变成 $0\cdot E=1$，无解不是概率模型矛盾，而是所求期望没有该模表示。

未约分表达式的分母不可逆，不代表最终答案一定不可表示。
表达式 $p/p=1$ 应先在整数或有理数层面约分，不能调用逆元求 $0/0$。
类似地，模线性系统奇异时，可能存在有理数层面的消去或约分，但不能从任意模特解恢复唯一的真实期望。
合数模数 $m$ 下，分母须满足 $\gcd(b,m)=1$，且不能用费马小定理求逆。

本库接口见第~\ref{compact-ModInt}~节（第~\pageref{compact-ModInt}~页）的 \texttt{ModInt<p>} 与
第~\ref{compact-GaussMod}~节（第~\pageref{compact-GaussMod}~页）的 \texttt{GaussMod<p>::solve(a,n)}。
前者的 \texttt{try\_inv()} 用欧几里得，允许合数模数的单位元，不可逆时返回空值；
\texttt{inv()} 与除法仍要求分母可逆。后者仍要求 $p$ 为素数，不能随基础类型扩展而放宽；
接收 $s$ 行 $s+1$ 列增广矩阵，可填为 $[I-Q\mid b]$，显式传入变量数 $s$。
先检查 \texttt{consistent}，再检查 \texttt{rank==s}，唯一解才读取 \texttt{particular}。
这些接口不会替你证明吸收、期望有限或有理答案可取模。

\section{SG 的前提、递推与异或策略}\label{knowledge-sg}\index{Sprague-Grundy}\index{mex}
先核对模型：双人轮流操作；双方在同一状态的合法操作完全相同（无偏游戏）；
没有随机性与隐藏信息；不能行动者输（正常玩法）；状态图为有限 DAG。
本节只在这些条件下使用普通 Sprague--Grundy 理论。“公平”在这里不表示胜率各半。

令 $g(u)$ 是状态的 SG 值，$\operatorname{mex}$ 表示集合中未出现的最小非负整数，则
\[
g(u)=\operatorname{mex}\{g(v):u\to v\},\qquad
\operatorname{mex}\varnothing=0.
\]
按逆拓扑顺序计算；若递归记忆化，栈深度仍可能等于最长路径长度。
若有 $d$ 个后继，则 $g(u)\le d$，只需标记 $0,\ldots,d$，超过 $d$ 的后继值可以忽略。
重复出现同一个后继值不影响 mex；用时间戳避免每个点都清空一张全局大表，
可做到 $O(|V|+|E|)$ 时间和 $O(|V|+|E|)$ 空间（包含存图）。

$g(u)=0$ 时所有后继非零；$g(u)>0$ 时一定存在后继值为零。
由 DAG 上归纳，零状态恰为先手必败态。不过 SG 数值不是胜率，也不是距离终点的步数。
若题目只问一个不可分解游戏的胜负，布尔递推“存在后继必败则必胜”已足够；
需要组合多个游戏时才必须保留比输赢更多的信息。

\textbf{异或只针对不交和。} 一步恰好选择一个子游戏走一步，其他子游戏的状态及合法操作不变，
且总游戏不能走时才输。此时组合值为 $G=g_1\oplus\cdots\oplus g_k$。
$G=0$ 时任意一步都把一个 $g_i$ 改为不同值，因此异或不再为零。
$G\ne0$ 时，取 $G$ 最高为一位上也为一的某个 $g_i$，令 $h=g_i\oplus G<g_i$。
mex 定义保证存在一步从该子状态走到 SG 值 $h$ 的后继，操作后整体异或为零。
实际输出方案时枚举这个子游戏的合法后继寻找 $h$，不能只输出一个目标数字。

如果一次允许同时改变两堆，或者一堆的操作会禁止另一堆的操作，就不是上述不交和。
例如两堆各有一颗石子，允许取走任意一堆或同时取走两堆：
把两堆当普通 Nim 会算出 $1\oplus1=0$，但先手同时取完立刻获胜。

\section{减法游戏小例与三类失效边界}\label{knowledge-subtraction-game}
一堆 $n$ 颗石子，每次可取 $S=\{1,3,4\}$ 中一种数量且不能取成负数。
空堆无合法操作，故 $g(0)=0$，对 $n>0$ 有
\[
g(n)=\operatorname{mex}\{g(n-s):s\in S,\ s\le n\}.
\]
从 $n=0$ 到 $13$ 逐项得到
\[
\begin{array}{c|rrrrrrr|rrrrrrr}
n&0&1&2&3&4&5&6&7&8&9&10&11&12&13\\\hline
g(n)&0&1&0&1&2&3&2&0&1&0&1&2&3&2
\end{array}
\]
例如 $g(4)=\operatorname{mex}\{g(3),g(1),g(0)\}=\operatorname{mex}\{1,0\}=2$，
不是合法操作数 3。两堆 $(4,5)$ 的异或为 $2\oplus3=1$，从第二堆取一颗得到 $(4,4)$ 即可归零。
两堆 $(1,1)$ 各自都是必胜态，但组合为必败态，不能把多个布尔“必胜”作逻辑或。

表中前两段相等本身不是周期证明。本例先验证 $g(n+7)=g(n)$ 对 $0\le n\le3$ 成立，
随后对 $n\ge4$，两边的后继都完整包含减 $1,3,4$；由更小下标的归纳假设，mex 集合相同，
故周期 7 对所有 $n\ge0$ 成立。对一般有限减法集合，SG 值不超过 $|S|$，
且从足够大的 $n$ 起只依赖前 $\max S$ 项，因此整个长度为 $\max S$ 的窗口重复后才保证后续周期。
它只保证最终周期，不能仅凭几项重复猜出从零开始的周期或短周期。

\begin{itemize}
\item \textbf{反常玩法（mis\`ere）。} 最后行动者输时，普通 SG 的终点规则与异或结论不可直接搬用。
只有一颗石子的 Nim 在正常玩法下先手赢，在反常玩法下先手输。
反常 Nim 有专门结论，不意味着所有减法游戏都能“把最终答案取反”。
\item \textbf{有环游戏（loopy）。} 单点自环将要求 $g=\operatorname{mex}\{g\}$，不存在这样的非负整数。
允许无限行走时还可能有和局；必须按题目规定处理重复局面、步数上限和无限对局。
若加入剩余步数后形成 DAG，可以在扩展状态上分析，不能把环当成已访问后直接返回零。
\item \textbf{有偏游戏（partizan）。} 双方在相同局面可走步不同，不满足无偏性。
仅一条“左方可走、右方不可走”的边，就不能用不带行动方信息的 mex 描述。
可按完整状态做胜负搜索，但一般不能将各部分的普通 SG 值异或。
\end{itemize}

\noindent\textbf{本地知识例证。}
\path{tests/math_knowledge_probability.py} 使用有理数精确运算检查指示变量、骰子尾和、
连续正面方程及模分母反例，并以直接胜负递归核对小规模减法游戏不交和。
它验证本节例子，不是通用概率求解器、SG 模板或在线评测 AC。

\label{end-probability-games-supplement}
